\subsection{Conditional residuals as the typing substrate}
The base detector~\citep{Wadinger2023} maintains, for each signal $i$, a conditional model of its current value $x_i$ given the remaining signals $x_{-i}$. We define the conditional residual
\begin{equation}
    r_i = x_i - \uis{\mu}{cond}{i}(x_{-i}),
\end{equation}
the difference between the observed value and its conditional mean, and the conditional standard deviation $\uis{\sigma}{cond}{i}$ already used by the detector to form its score. Both quantities are available at every step at no additional modelling cost. For brevity we write $\sigma$ for the per-channel conditional standard deviation $\uis{\sigma}{cond}{i}$ in the results, and we also use $\sigma$ as a \emph{magnitude unit}---a fault or shift ``of $k\sigma$'' means $k$ multiples of the relevant conditional standard deviation; the intended sense is clear from context. The classifier operates entirely on the stream of residuals $\{r_i\}$ and their conditional scale $\uis{\sigma}{cond}{i}$: it never stores past samples and never refits a model over a window. This keeps the per-signal state to a fixed set of scalars and the per-step cost constant, matching the streaming, O(1)-memory footprint of the detector it extends.

\subsection{Running statistics}
For each signal the classifier maintains a small set of exponentially weighted moving (EWM) statistics over the residual stream. We track a fast (short half-life) and a slow (long half-life) EWM mean of $r_i$, denoted $\uis{m}{fast}{i}$ and $\uis{m}{slow}{i}$, and an EWM variance $\uis{v}{}{i}$ of $r_i$. We additionally track the EWM variance of the raw signal increments to detect a stalled transmitter. All updates are the standard constant-memory recursions; no buffer of past values is required. The half-lives are the only temporal hyperparameters of the typing layer and are shared across signals unless per-signal calibration is applied (Section~\ref{Discussion}).

\subsection{The four type tests}
Given the running statistics, a flagged signal is attributed to a type by the following tests, each expressed in units of the conditional scale so that thresholds are dimensionless.

\paragraph{Bias.} A persistent offset of the residual from zero, with stable variance, indicates a constant calibration error. We declare \emph{bias} when the slow residual mean is large in conditional-scale units, $|\uis{m}{slow}{i}| / \uis{\sigma}{cond}{i}$ exceeding a threshold, while the residual variance remains close to its baseline.

\paragraph{Drift.} A gradually growing offset is distinguished from a static one by the \emph{gap between the short- and long-horizon EWM means}. We declare \emph{drift} when $|\uis{m}{fast}{i} - \uis{m}{slow}{i}| / \uis{\sigma}{cond}{i}$ exceeds a threshold, i.e.\ when the recent residual mean is moving away from its longer-run value. We name this honestly: it is a short-versus-long EWM mean-gap statistic, \emph{not} a fitted regression slope. It requires no stored window and no least-squares fit; it responds to the direction and recency of residual movement rather than estimating a rate.

\paragraph{Accuracy loss.} An inflated noise floor with an unchanged mean indicates degraded measurement accuracy. We declare \emph{accuracy loss} when the residual variance $\uis{v}{}{i}$ exceeds its baseline by a multiplicative factor while the slow residual mean stays near zero.

\paragraph{Freezing.} A transmitter stuck at its last value produces near-zero signal increments while the conditional expectation continues to move. We declare \emph{freezing} when the EWM variance of the raw increments collapses below an absolute floor $\texttt{freeze\_eps}$ and below a fraction $\texttt{freeze\_var\_ratio}$ of the signal's marginal variance, scaled by $\texttt{freeze\_abs\_scale}$. These three constants govern how aggressively a quiet signal is judged frozen; their calibration is examined in Section~\ref{Results} (the freeze-on-quiescent ablation). A signal that satisfies none of the tests while flagged is labelled \emph{normal} with respect to typing (e.g.\ a peer signal pulled by cross-talk), and an unflagged signal is \emph{normal} by default.

\subsection{Change-point graded suppression}
A fault that arrives while the base detector is adapting to a genuine regime shift can be either masked by the adaptation or mistaken for one. The detector's change-point machinery applies a graded suppression to the typing decision around detected change points, controlled by a scale parameter \texttt{suppress\_scale}: at \texttt{scale}~$=1$ no suppression is applied, the default \texttt{scale}~$=5$ applies graded suppression, and \texttt{scale}~$=\infty$ suppresses any typing within the change-point window entirely. This parameter trades off two failure modes that we quantify separately in Section~\ref{Results}: spurious type labels from coordinated shifts in weakly correlated signals (reduced by stronger suppression) against a masking dead-band for genuine faults co-occurring with a change point (widened by stronger suppression). We treat \texttt{suppress\_scale} as the principal tuning knob of the composed system and report its effect rather than fixing it silently.
